\section{Related Work}
\label{sec:related}

\para{Memorization in neural networks.}
\citet{carlini2023quantifying} establish three log-linear relationships governing memorization in language models: memorization grows with model capacity, data duplication, and context length. \citet{zhang2023counterfactual} refine this picture by distinguishing extractable memorization from counterfactual memorization, showing that many apparently memorized sequences are simply predictable from other data. \citet{feldman2020does} argue that some memorization is \emph{necessary} for good generalization on long-tailed distributions, since rare subpopulations must be memorized to avoid systematic errors. Our work does not dispute this---rather, we argue that the \emph{speed} of memorization, not memorization itself, is the bottleneck.

\para{Forgetting dynamics.}
\citet{toneva2019empirical} introduce forgetting events---transitions from correct to incorrect classification during training---and show that examples vary dramatically in forgettability. Roughly 30\% of CIFAR-10 examples are ``unforgettable'' (learned once, never forgotten), and removing them preserves test accuracy while enabling 30\% dataset compression. \citet{zheng2025spurious} distinguish spurious forgetting (loss of task alignment) from true knowledge loss in continual learning. We adopt Toneva et al.'s forgetting event framework and extend it to the grokking setting, tracking how forgetting pressure modulates per-example learning dynamics.

\para{Grokking and delayed generalization.}
\citet{power2022grokking} discover that small transformers trained on modular arithmetic achieve perfect training accuracy thousands of steps before test accuracy improves. Weight decay is critical for this transition. \citet{li2025grokking} demonstrate that grokking is not limited to toy settings: using OLMoE-1B-7B, they observe a clear temporal gap between memorization and generalization onset during LLM pretraining, detectable via routing pathway entropy. Unlike these works, which \emph{observe} the memorize-then-generalize transition, we systematically \emph{control} it by varying forgetting pressure and measuring the causal effect on generalization speed.

\para{Compression and representation learning.}
\citet{yu2025memorization} show that LLM pretraining alternates between memorization phases (expanding representations) and compression phases (contracting them). Adding an explicit compression loss improves MMLU by 1.8\% and HellaSwag by 2.2\%. \citet{wang2025breaking} identify a ``memorization barrier'' in code fine-tuning---models pre-memorize downstream data during pretraining, creating bad local optima---and break through it using an Information Bottleneck approach. Both works support our hypothesis that compression (a form of forgetting) enables generalization. We complement these findings with controlled experiments that isolate forgetting pressure as the key variable.

\para{Data deduplication and efficiency.}
\citet{lee2022deduplicating} show that standard datasets contain massive duplication, and that deduplication reduces memorized output by 10$\times$ while improving training efficiency. \citet{abbas2023semdedup} extend this to semantic deduplication, removing 50\% of data with minimal in-distribution performance loss and improved out-of-distribution generalization. \citet{muennighoff2023scaling} find that up to 4 epochs of repetition approximates unique data, but returns diminish rapidly beyond that. \citet{tirumala2023d4} show that intelligent data selection via embeddings provides 20\% efficiency gains. Our Experiment 2 directly tests the effect of controlled duplication on grokking dynamics, complementing these large-scale deduplication studies with a mechanistic perspective.
